\section{Étape 2 -- Configuration matérielle du coupleur}
\label{secEtapeConfigMaterielle}

Cette étape consiste à déclarer le coupleur BMX~NOC~0401.2 dans la
configuration matérielle du rack de l'automate, à l'aide de l'éditeur
graphique de configuration (double clic sur le dossier \emph{Configuration}
de l'explorateur du projet).

\subsection{Déclaration de l'emplacement}

Dans l'emplacement (slot) du rack destiné à recevoir le coupleur, ajouter
un module \textbf{BMX~NOC~0401.2} (module 4 ports Ethernet RJ45 10/100).
Sur la cellule considérée, deux coupleurs sont ainsi déclarés, un par
emplacement.

\subsection{Renseignement du DTM}

L'ajout d'un module BMX~NOC~0401.2 fait apparaître une boîte de dialogue à
plusieurs onglets, dans laquelle il convient de renseigner :

\begin{itemize}
    \item l'\textbf{onglet Général} : le nom du DTM à donner au coupleur
    (par exemple \texttt{NOC\_ETHIP\_CONVOYEUR} ou
    \texttt{NOC\_ETHIP\_ROBOT}) ;
    \item l'\textbf{onglet Information sur le protocole} : le protocole à
    utiliser, à savoir \textbf{CIP (EtherNet/IP)} -- il s'agit normalement
    du protocole sélectionné par défaut.
\end{itemize}

\subsection{Configuration complémentaire du processeur}

Si l'application n'utilise pas le réseau CANopen intégré au processeur, il
convient de configurer sa voie en ne réservant aucun objet \texttt{\%M} ni
\texttt{\%MW} pour ce réseau, ni en entrée ni en sortie (double clic sur la
prise DB9 du réseau CANopen du processeur pour faire apparaître sa boîte
de dialogue de configuration).

\subsection{Sauvegarde et vérification}

Une fois la configuration matérielle établie :

\begin{enumerate}
    \item sauvegarder l'application ;
    \item la compiler ;
    \item vérifier, à l'aide de l'éditeur de données (variables dérivées),
    que les zones mémoire allouées aux coupleurs BMX~NOC~0401.2
    correspondent à l'attendu. La compilation crée automatiquement, pour
    chaque coupleur, les types de données nécessaires pour représenter ses
    zones d'entrée et de sortie.
\end{enumerate}
